\IfFormatAtLeastTF{2025-06-01}{\DocumentMetadata{lang=en-US,pdfversion=2.0,tagging=on}}{}
\documentclass[11pt]{amsart}
\usepackage[T1]{fontenc}
\IfFileExists{lmodern.sty}{\usepackage{lmodern}}{\PassOptionsToPackage{expansion=false}{microtype}}
\usepackage{amsmath,amssymb,amsthm,mathtools}
\usepackage{etoolbox}
\usepackage[letterpaper,textwidth=5.6in,textheight=\textheight,hmarginratio=1:1]{geometry}
\usepackage[nopatch=footnote]{microtype}
\usepackage{xcolor}
\usepackage[colorlinks=true,linkcolor=blue,citecolor=blue,urlcolor=blue]{hyperref}
\newif\ifanonymous
\anonymousfalse
\ifanonymous
  \newcommand{\paperauthor}{Anonymous}
\else
  % Repository locator for the public (non-anonymous) build of
% non_mf_groups_exist.tex.  Kept in a separate file so that the anonymous
% build has an anonymous source archive and not merely an anonymous PDF:
% exclude this file from a double-blind submission and the manuscript still
% compiles, printing "supplied with this submission as supplementary
% material" in place of the link.
%
% Loaded only when \anonymousfalse is set in the preamble.
\newcommand{\coderepository}{https://github.com/SauersML/group-approximation}
\newcommand{\coderepositoryname}{github.com/SauersML/group-approximation}

  \newcommand{\paperauthor}{\paperauthorname}
\fi
\hypersetup{pdftitle={A simple unital C*-algebra that is not K_1-injective},pdfauthor={\paperauthor},pdfsubject={K_1-injectivity, simple C*-algebras, Problem LIX of Schafhauser-Tikuisis-White},pdfkeywords={K_1-injectivity, simple C*-algebras, AH algebras, Murray-von Neumann equivalence, Wu formula, Steenrod squares}}
\newcommand{\doi}[1]{\href{https://doi.org/#1}{\nolinkurl{doi:#1}}}
% Machine-readable source metadata attaching a printed statement to an
% accompanying Lean~4 declaration.  The macro is intentionally invisible in
% the PDF; the web edition and audit scripts consume its two arguments.
% Each \leanverified marker identifies a selected printed statement and a Lean
% declaration intended to prove it.  The markers do not assert line-by-line
% verification of the prose or complete sentence-census coverage of the
% manuscript.
\newcommand{\leanverified}[2]{\ignorespaces}
\allowdisplaybreaks
\raggedbottom
\emergencystretch=1.5em
\newtheorem{theorem}{Theorem}[section]
\newtheorem{proposition}[theorem]{Proposition}
\newtheorem{lemma}[theorem]{Lemma}
\newtheorem{corollary}[theorem]{Corollary}
\theoremstyle{plain}
\newtheorem{mainthm}{Theorem}
\renewcommand{\themainthm}{\Alph{mainthm}}
\theoremstyle{definition}
\newtheorem{remark}[theorem]{Remark}
\newcommand{\F}{\mathbb F}
\newcommand{\C}{\mathbb C}
\newcommand{\R}{\mathbb R}
\newcommand{\CP}{\mathbb{CP}}
\newcommand{\U}{\operatorname{U}}
\newcommand{\Uz}{\operatorname{U}_0}
\newcommand{\Sq}{\operatorname{Sq}}
\newcommand{\rank}{\operatorname{rank}}
\newcommand{\tr}{\operatorname{tr}}

\title{A simple unital C*-algebra that is not $K_1$-injective}
\author{\paperauthor}
\date{\today}

\begin{document}
\begin{abstract}
We answer a question that Blackadar's book records as open and that
Schafhauser, Tikuisis and White pose as Problem LIX of their list: there is a
unital simple separable C*-algebra $A$ and a unitary $v \in \U(A)$ with
$v \notin \Uz(A)$ and $\operatorname{diag}(v,1) \in \Uz(M_2(A))$, so the
canonical map $\U(A)/\Uz(A) \to K_1(A)$ is not injective.  The algebra is an
inductive limit of homogeneous algebras over products of a sphere with complex
projective spaces, with twisted point evaluations as connecting maps.  The
unitary $v$ comes from the nonzero element of $\pi_4(\U(2)) \cong \mathbb Z/2$,
and the proof that $v$ is not null-homotopic in $A$ reduces to the
non-isomorphism of two vector bundles of the same rank, which we prove by
computing the top mod-2 Chern class of a mapping torus in two ways.  The proof
is formalised in Lean~4.
\end{abstract}

\maketitle

\section{Introduction}
\label{sec:intro}

A unital C*-algebra $A$ is \emph{$K_1$-injective} when the canonical map
\[
  \kappa_A \colon \U(A)/\Uz(A) \longrightarrow K_1(A)
\]
is injective, where $\Uz(A)$ is the connected component of the identity in the
unitary group $\U(A)$.  Since $K_1(A)$ is the direct limit of the groups
$\U(M_n(A))/\Uz(M_n(A))$ along $u \mapsto \operatorname{diag}(u,1)$, the map
$\kappa_A$ is injective when every unitary $v \in \U(A)$ with
$\operatorname{diag}(v,1,\dots,1) \in \Uz(M_n(A))$ for some $n$ already lies in
$\Uz(A)$.

Injectivity fails for some nonsimple algebras.  Blackadar
\cite[8.1.2(d)]{Bl} gives the classical example, which rests on a theorem of
Araki, James and Thomas \cite{AJT}: for $A = C(\U(2) \times \U(2), M_2)$ and
the two coordinate unitaries $u$ and $v$, the products $uv$ and $vu$ lie in
different components of $\U(A)$, while $\operatorname{diag}(uv,1)$ and
$\operatorname{diag}(vu,1)$ lie in the same component of $\U(M_2(A))$.  For
simple algebras the map is injective in every regularity class where it has
been computed: purely infinite simple algebras \cite{Cu}, algebras of stable
rank one \cite{Ri}, algebras of real rank zero \cite{Li}, and
$\mathcal Z$-stable algebras \cite{Ji}.  Blackadar wrote in 1998 that for
simple unital $A$ the map is expected to be an isomorphism, ``although this
remains an open problem'' \cite[\S 8.1]{Bl}.  Villadsen \cite{Vi02} then showed
that surjectivity can fail for simple separable unital nuclear algebras.
Schafhauser, Tikuisis and White record injectivity as the open direction and
pose it as Problem LIX of their list \cite{STW}: is every unital simple
C*-algebra $K_1$-injective?  The answer is no.

\begin{mainthm}
\label{thm:main}
There is a unital simple separable C*-algebra $A$ that is not
$K_1$-injective.  Explicitly, there is a unitary $v \in \U(A)$ with
$v \notin \Uz(A)$ and $\operatorname{diag}(v,1) \in \Uz(M_2(A))$.
\leanverified{thm:main}{GroupApproximation.NinetyNineProblems.not\_problemLIX}
\leanverified{thm:main-separable}{GroupApproximation.NinetyNineProblems.exists\_separable\_simple\_unital\_not\_k1Inj}
\end{mainthm}

Simplicity means that every closed two-sided ideal of $A$ is $0$ or $A$; for
a unital C*-algebra this is the same as simplicity as a ring.

The algebra $A$ is an inductive limit of homogeneous algebras
$A_j = E_j\, M(C(X_j))\, E_j$, where $X_j = S^4 \times \prod_{i<j} \CP^{d_i}$
and $E_j$ is the direct sum of a trivial rank-two bundle with multiples of the
tautological line bundles of the projective factors.  The connecting maps are
twisted point evaluations $a \mapsto a \oplus a(x_j) \otimes 1_{L}$, with $L$
the tautological line bundle of the new factor.  Inductive limits of this kind
were introduced by Villadsen \cite{Vi98,Vi99}.  Point evaluations at a sequence
of points with dense tails make the limit simple, and the twist by $L$ is
necessary: a simple inductive limit of this shape with untwisted point
evaluations has stable rank one \cite{EHT}, so is $K_1$-injective by
\cite{Ri}, and is even $K$-stable \cite{Se}.

The unitary $v$ is the image in $A$ of the clutching function
$u \colon S^4 \to \U(2)$ of the rank-two bundle $F = x^{\perp}$ over
$S^5 \subset \C^3$; it represents the nonzero element of
$\pi_4(\U(2)) \cong \mathbb Z/2$ \cite{Bo}.  Since $\operatorname{diag}(u,1)$
is null-homotopic in $\U(3)$, $\operatorname{diag}(v,1) \in \Uz(M_2(A))$.  The
content of the theorem is $v \notin \Uz(A)$.  A null-homotopy of $v$ in $A$
gives one at a finite stage $A_j$, and a null-homotopy at stage $j$ gives an
isomorphism of vector bundles $F \oplus H \cong \mathbf 1^2 \oplus H$ over
$S^5 \times \prod_{i<j} \CP^{d_i}$, where $H$ is the sum of the twisted line
bundles.  Lemma~\ref{lem:two} says that no such isomorphism exists.  Its proof
is a mod-2 computation: from an isomorphism one builds a bundle of rank $r$
over the manifold $S^1 \times S^5 \times \prod_{i<j} \CP^{d_i}$ of dimension
$2r$, whose top mod-2 Chern class is nonzero because the bundle has a section
with a single nondegenerate zero, and zero by the Wu formula because the
dimensions $d_i$ are even.

The whole proof is formalised in Lean~4 over Mathlib, in the repository
\href{https://github.com/SauersML/group-approximation}{SauersML/group-approximation};
the argument given here is the one formalised.  Independently, Toms \cite{To}
has given a different solution of Problem LIX, starting from the same class in
$\pi_4(\U(2))$ but trapping it with higher-rank matrices and a spin bordism
obstruction.

\section{The algebra}
\label{sec:algebra}

\subsection{Bundles as projections}

Write $\CP^d$ for the space of rank-one projections in $M_{d+1}(\C)$, $S^4$
for the four-sphere and $S^5$ for the unit sphere of $\C^3$.  A complex vector
bundle over a compact Hausdorff space $X$ is a continuous projection-valued
map, that is, a projection $p \in M_K(C(X))$; an isomorphism of bundles is a
Murray--von Neumann equivalence of projections, and the algebra of sections of
the endomorphism bundle of $p$ is the corner $p\, M_K(C(X))\, p$.

\subsection{The tower}

Set $d_i = 2^{\,i+1}$ and, for $j \geq 0$,
\[
  Y_j \;=\; \prod_{i<j} \CP^{d_i},
  \qquad
  X_j \;=\; S^4 \times Y_j,
  \qquad
  H_j \;=\; \bigoplus_{i<j} L_i^{\oplus d_i},
\]
where $L_i$ is the tautological line bundle of the $i$-th projective factor,
pulled back to $Y_j$.  Put $E_j = \mathbf 1^2 \oplus H_j$ and
\[
  A_j \;=\; E_j \, M(C(X_j)) \, E_j .
\]
The bundle $E_j$ has rank $2 + \sum_{i<j} d_i = 2^{\,j+1}$ and the base $X_j$
has dimension $4 + 2\sum_{i<j} d_i = 2^{\,j+2}$, twice the rank.  The
dimensions $d_i$ are even; Step~D in Section~\ref{sec:modtwo} uses this.

\subsection{The connecting maps and the limit}

Since $X_{j+1} = X_j \times \CP^{d_j}$ and $E_{j+1} = E_j \oplus L_j^{\oplus
d_j}$, and since the multiplicity $d_j$ equals the rank of $E_j$, an element
$a(x_j)$ of the fibre of $E_j M E_j$ at a point $x_j \in X_j$ acts on
$L_j \otimes \C^{d_j} = L_j^{\oplus d_j}$.  The connecting map
$\varphi_j \colon A_j \to A_{j+1}$ is
\[
  \varphi_j(a)(x,z) \;=\; a(x) \,\oplus\, \bigl(a(x_j) \otimes 1_{L_j,\,z}\bigr),
  \qquad x \in X_j,\ z \in \CP^{d_j},
\]
where the base points $x_j \in X_j$ are chosen so that, for every $k$, the
images of $x_j$ for $j \geq k$ under the coordinate projections $X_j \to X_k$
are dense in $X_k$.  Each $\varphi_j$ is an injective unital $*$-homomorphism.
Let $A$ be the inductive limit, with injective unital maps
$\iota_j \colon A_j \to A$ whose images have dense union.  The algebra $A$ is
separable, since each $A_j$ is.

$A$ is simple.  Let $0 \neq a \geq 0$ in $A_k$ and pick $i \geq k$ such that
the image of $x_i$ in $X_k$ lies in the open set where $a$ does not vanish.
Then the point-evaluation summand of $\varphi_{k,i+1}(a)$ is nonzero at every
point of $X_{i+1}$.  A section that is nonzero in every fibre generates the
corner as a closed ideal, because the fibres are full matrix algebras.  So
every nonzero closed two-sided ideal of $A$ contains $1$.

\subsection{The generator}

Over $S^5$, the projection $F(x) = \mathbf 1_3 - x x^{*}$ has rank two.  For
each pole $p = \pm e_3$, two Householder reflections through the normalised
midpoint of $p$ and $x$ give a continuous unitary-valued map $\sigma_p$ on the
complement of the antipode of $p$ with $\sigma_p(x)\, e_3 = x$, so
$\sigma_p(x)$ maps $e_3^{\perp}$ onto $F(x)$.  On the equatorial $S^4$ the two
frames differ by the unitary
\[
  u \colon S^4 \longrightarrow \U(2), \qquad u(x) = \sigma_+(x)^{*}\sigma_-(x),
\]
which fixes $e_3$ and so takes values in the unitary group of $e_3^{\perp}$.
Thus $u$ is the clutching function of $F$: $F$ is obtained from the trivial
rank-two bundles $\sigma_{\pm}(e_3^{\perp})$ over the two closed hemispheres by
gluing along $u$.  Since $\sigma_+$ extends over the northern hemisphere and
$\sigma_-$ over the southern one, $\operatorname{diag}(u,1) = \sigma_+^{*}
\sigma_-$ is null-homotopic in $\U(3)$.

At stage $0$ we have $H_0 = 0$ and $A_0 = M_2(C(S^4))$, so $u \in \U(A_0)$.
The witness of Theorem~\ref{thm:main} is $v = \iota_0(u)$.

\section{Reduction to a statement about bundles}
\label{sec:reduction}

Write $u_j$ for the image of $u$ in $A_j$.  Each connecting map adds a block
evaluated at a fixed point, so
\[
  u_j \;=\; u \oplus c_j
\]
for a unitary $c_j$ of $H_j$ over $Y_j$; in particular $c_j$ does not depend
on the $S^4$ coordinate.

\begin{lemma}
\label{lem:stage}
If $u_j \in \Uz(A_j)$, then $F \oplus H_j \cong \mathbf 1^2 \oplus H_j$ over
$S^5 \times Y_j$.
\end{lemma}

\begin{proof}
A unitary $g$ of $\mathbf 1^2 \oplus H_j$ over $S^4 \times Y_j$ determines a
bundle over $S^5 \times Y_j$, obtained by gluing the restrictions of
$\mathbf 1^2 \oplus H_j$ to the two closed hemispheres along $g$.  Homotopic
unitaries give isomorphic bundles, a unitary that extends over one hemisphere
gives $\mathbf 1^2 \oplus H_j$, and gluing along $u \oplus 1$ gives
$F \oplus H_j$.  Since $c_j$ does not depend on the sphere coordinate,
$1 \oplus c_j$ extends over a hemisphere, so gluing along
$u \oplus c_j = (u \oplus 1)(1 \oplus c_j)$ still gives $F \oplus H_j$.  If
$u \oplus c_j$ is null-homotopic, this bundle is isomorphic to
$\mathbf 1^2 \oplus H_j$.
\end{proof}

\begin{lemma}
\label{lem:limit}
If $v \in \Uz(A)$, then $u_j \in \Uz(A_j)$ for some $j$.
\end{lemma}

\begin{proof}
A path from $v$ to $1$ in $\U(A)$ is compact, so for some $j$ it lies within
distance $\tfrac12$ of $\iota_j(A_j)$.  Choose finitely many points along the
path, consecutive ones within distance $\tfrac12$ of each other, and
approximate each by an element of $A_j$ within $\tfrac12$.  Each approximant
is invertible and close to its unitary part, and consecutive unitary parts are
at distance less than $2$, so the spectrum of the product of one with the
adjoint of the next avoids $-1$ and the functional calculus joins them by a
path in $\U(A_j)$.  The path starts at $u_j$ and ends at $1$.
\end{proof}

\begin{proof}[Proof of Theorem~\ref{thm:main} from Lemma~\ref{lem:two}]
If $v \in \Uz(A)$, Lemma~\ref{lem:limit} gives $u_j \in \Uz(A_j)$ for some
$j$, and Lemma~\ref{lem:stage} gives $F \oplus H_j \cong \mathbf 1^2 \oplus
H_j$ over $S^5 \times Y_j$, contradicting Lemma~\ref{lem:two} with $l = j$.
So $v \notin \Uz(A)$.  On the other hand
$\operatorname{diag}(u,1) \in \Uz(M_2(A_0))$, since it is null-homotopic
already in $\U(3)$, and applying $\iota_0 \otimes \mathrm{id}_{M_2}$ gives
$\operatorname{diag}(v,1) \in \Uz(M_2(A))$.
\end{proof}

\section{The bundle obstruction}
\label{sec:modtwo}

\begin{lemma}
\label{lem:two}
Let $d_0,\dots,d_{l-1}$ be even and positive, and let
\[
  M \;=\; S^5 \times \prod_{i<l} \CP^{d_i},
  \qquad m = \sum_{i<l} d_i ,
  \qquad H = \bigoplus_{i<l} L_i^{\oplus d_i}.
\]
Over $M$ the bundles $F \oplus H$ and $\mathbf 1^2 \oplus H$, both of rank
$2+m$, are not isomorphic.
\leanverified{lem:two}{GroupApproximation.LIX.LemmaTwoHolds}
\end{lemma}

Both bundles are subbundles of $V = \mathbf 1^3 \oplus H$, of rank $r = 3+m$:
with $e$ the constant section $e_3$ of $\mathbf 1^3$ and $s$ the section
$x \mapsto x$, we have $F \oplus H = V - ss^{*}$ and
$\mathbf 1^2 \oplus H = V - ee^{*}$.  Since the two bundles have the same rank,
no dimension count separates them.

The argument is a mod-2 characteristic class computation.  We write
$\gamma_i$ for the mod-2 Chern classes and $\gamma = \sum_i \gamma_i$ for the
total class, and all cohomology has coefficients in $\F_2$.  The source
manuscript proves the lemma integrally, by a Chern number computation with a
K\"unneth decomposition over $\mathbb Z$; the argument below, which is the one
formalised, works mod~$2$ throughout.  Working mod~$2$ removes signs and
torsion, replaces the K\"unneth theorem over $\mathbb Z$ by one over a field,
and makes the evenness of the $d_i$ enter through the Frobenius endomorphism
of an $\F_2$-algebra.

\subsection{The setting}

Fix even positive $d_0,\dots,d_{l-1}$, put $m = \sum_i d_i$ and $r = 3+m$, and
let $M$, $H$ and $V$ be as above.  Everything happens over the closed manifold
\[
  N \;=\; S^1 \times M
\]
of dimension $1 + 5 + 2m = 2r$.  So $H^{2r}(N)$ is the top cohomology group of
$N$, it is one-dimensional over $\F_2$, and it is the only group in which a
rank-$r$ bundle over $N$ can have a nonzero top Chern class.  We compute
$\gamma_r$ of one bundle over $N$ in two ways and get different answers.

\subsection{Step A: from an isomorphism to a unitary}

Suppose $F \oplus H \cong \mathbf 1^2 \oplus H$, that is, $V - ss^{*}$ and
$V - ee^{*}$ are Murray--von Neumann equivalent.  Extending an equivalence by
$e \mapsto s$ gives a continuous unitary $g$ of the corner of $V$ with
$g e = s$.

\subsection{Step B: the mapping torus}

Cut $S^1$ at two points into two arcs and let $W_g$ be the bundle over
$N = S^1 \times M$ obtained from $V \times \mathrm{arc}$ on each arc by gluing
with the identity at one seam and with $g$ at the other.  Then $W_g$ has rank
$r$ and restricts to $V$ on each slice $\{t\} \times M$.

The rotation model, in which one runs a path from the identity to $g$ around
the circle, does not work here: the rotation returns to the identity, so the
bundle it produces is trivial.

\subsection{Step C: $\gamma_r(W_g) \neq 0$}

The bundle $V$ has the section $e$, and $H$ has a continuous section with a
single nondegenerate zero, obtained by choosing, on each factor $\CP^{d_i}$,
$d_i$ sections of $L_i$ that vanish simultaneously at one point only.
Interpolating between $e$ and $s$ along the arc, and switching the section of
$H$ on in the middle of it, gives a section $\sigma$ of $W_g$ with exactly one
zero $z \in N$.  The zero is nondegenerate: in a chart at $z$, the section is a
local homeomorphism onto a neighbourhood of the origin in $\C^r$, by the
inverse function theorem.

Let $U \in H^{2r}(W_g, W_g \setminus 0)$ be the Thom class of $W_g$.  Pulling
back along $\sigma$, viewed as a map of pairs $(N, N \setminus z) \to
(W_g, W_g \setminus 0)$, gives $\sigma^{*}U \in H^{2r}(N, N \setminus z)$, and
by the Gysin identity of Lemma~\ref{lem:hclass} the image of $\sigma^{*}U$
in $H^{2r}(N)$ is $\sigma^{*}\pi^{*}\gamma_r(W_g) = \gamma_r(W_g)$.  Since
$\sigma$ has a single nondegenerate zero, $\sigma^{*}U$ is the generator of
$H^{2r}(N, N \setminus z) \cong \F_2$ (Lemmas~\ref{lem:hres} and
\ref{lem:hsq}), and the map $H^{2r}(N, N \setminus z) \to H^{2r}(N)$ is an
isomorphism, by a Mayer--Vietoris computation for the punctured manifold
$N \setminus z$.  So $\gamma_r(W_g) \neq 0$.

\subsection{Step D: $\gamma_r(W) = 0$}

Let $W$ be any rank-$r$ bundle over $N$ that restricts to $V$ on a slice, and
let $p \colon N \to M$ be the projection.  By the K\"unneth theorem over
$\F_2$,
\[
  H^{*}(N) \;=\; \Lambda(t) \otimes \Lambda(x) \otimes H^{*}(Y),
\]
with $t$ the class of $S^1$, $x$ the class of $S^5$ and
$Y = \prod_i \CP^{d_i}$.  Because $W$ restricts to $V$ on a slice,
\[
  \gamma(W) \;=\; \gamma(p^{*}V)\,(1 + z\,b),
  \qquad z = t x,\quad b = \sum_n b_n \in H^{*}(Y),
\]
with $b_n \in H^{2n}(Y)$.  Three facts finish the computation.

First, $\gamma(V) = \prod_i (1+h_i)^{d_i}$, where $h_i \in H^2(\CP^{d_i})$ is
the mod-2 Chern class of $L_i$.  Since each $d_i$ is even and
$(1+h)^2 = 1 + h^2$ over $\F_2$, this is a polynomial in the $h_i^2$, so
$\gamma_k(V) = 0$ for every odd $k$.  This is the only place where the
evenness of the $d_i$ is used.

Second, the Wu formula for mod-2 Chern classes \cite{MS},
\[
  \Sq^{2i}\gamma_{i+1} \;=\; \sum_{s=0}^{i} \gamma_{i-s}\,\gamma_{i+1+s},
\]
applied to $\gamma(W)$ and read on the $z$-part, together with the
instability of the Steenrod squares ($\Sq^{2k}$ vanishes on classes of degree
below $2k$), gives $b_n = 0$ for every even $n$.

Third, $\gamma(V)$ lives in degrees at most $2m < 2r$, so the degree-$2r$
part of $\gamma(W)$ is
\[
  \gamma_r(W) \;=\; z \sum_{k} \gamma_{m-k}(V)\, b_k .
\]
Since $m$ is even, $\gamma_{m-k}(V) \neq 0$ forces $k$ even, and then
$b_k = 0$.  So $\gamma_r(W) = 0$.

\subsection{Conclusion}

Step~D applies to $W_g$, which restricts to $V$ on a slice, and gives
$\gamma_r(W_g) = 0$, while Step~C gives $\gamma_r(W_g) \neq 0$.  So no
unitary $g$ as in Step~A exists, and $F \oplus H$ and $\mathbf 1^2 \oplus H$
are not isomorphic.  This proves Lemma~\ref{lem:two}.
\leanverified{lem:two}{GroupApproximation.CharClass.lemmaTwo\_of\_stepC\_stepD}

\section{Three lemmas used in Step C}
\label{sec:stepc}

Throughout this section $E = W_g$ is the rank-$r$ bundle of Step~B over the
$2r$-manifold $N$, $\pi \colon E \to N$ is the projection, $U$ is its Thom
class, $\sigma$ is the section of Step~C and $z$ is its zero.  The Thom class
is constructed on the projective completion $P(E \oplus \mathbf 1)$: by the
Leray--Hirsch theorem, $H^{*}(P(E \oplus \mathbf 1))$ is a free module over
$H^{*}(N)$ with basis $1, \xi, \dots, \xi^{r}$, where $\xi$ is the mod-2 Chern
class of the tautological line bundle, and $U$ is the class whose
Leray--Hirsch coordinates are $(\gamma_r, \gamma_{r-1}, \dots, \gamma_1, 1)$.
The total space $E$ sits in $P(E \oplus \mathbf 1)$ as the affine chart, the
complement of $P(E)$, and $\xi$ restricts to zero on it.

\begin{lemma}[Gysin identity]
\label{lem:hclass}
The image of $U$ under $H^{2r}(E, E \setminus 0) \to H^{2r}(E)$ is
$\pi^{*}\gamma_r(E)$.
\leanverified{sec:hclass}{GroupApproximation.CharClass.lixHclass}
\end{lemma}

\begin{proof}
The image of $U$ in $H^{2r}(P(E \oplus \mathbf 1))$ is
$\sum_{i=0}^{r} \pi^{*}(\gamma_{r-i}) \smile \xi^{i}$.  Restricting to the
affine chart $E$ kills $\xi$, so only the term $i = 0$ survives, and it is
$\pi^{*}\gamma_r$.
\end{proof}

\begin{lemma}
\label{lem:hres}
Let $B$ be a trivialising open neighbourhood of $z$ in $N$.  The restriction
\[
  H^{2r}(E,\,E \setminus 0) \longrightarrow
  H^{2r}\bigl(E|_B,\,E|_B \setminus 0\bigr)
\]
is injective.
\leanverified{sec:hres}{GroupApproximation.CharClass.injective\_lixRes}
\end{lemma}

\begin{proof}
The Thom isomorphism identifies $H^{2r}(E, E \setminus 0)$ with $H^0(N)$,
which is one-dimensional because $N$ is connected, with generator $U$.  So it
suffices to show that $U$ restricts to a nonzero class over $B$.  Since
$z \in B$, restriction to the fibre over $z$ factors through restriction over
$B$, so it suffices that $U$ restricts to a nonzero class over the point $z$.
Leray--Hirsch coordinates restrict coordinatewise, and the last coordinate of
$U$ is $1 \in H^0(N)$, which restricts to $1 \in H^0(\{z\}) \neq 0$.
\end{proof}

\begin{lemma}
\label{lem:hsq}
Let $\phi$ be a chart of $N$ at $z$ in which $\sigma$ is a local
homeomorphism onto a neighbourhood of the origin in $\C^r$.  The class
$\sigma^{*}U \in H^{2r}(N, N \setminus z)$, read in the chart $\phi$, agrees
with the restriction of $U$ to the fibre $(E_z, E_z \setminus 0)$, read in a
trivialisation of $E$ near $z$.  In particular $\sigma^{*}U$ generates
$H^{2r}(N, N \setminus z) \cong \F_2$.
\leanverified{sec:hsq}{GroupApproximation.CharClass.lixHsq}
\end{lemma}

\begin{proof}
Both sides are pullbacks of $U$ along maps of pairs into
$(E, E \setminus 0)$, but reading $\sigma^{*}U$ in the chart involves the
inverse of an excision isomorphism, so the two sides cannot be compared
pointwise.  Composing both with the excision itself removes the inverse.  We
then restrict to a small neighbourhood of $z$, on which the chart reads the
section; only injectivity of this restriction is used.  What remains are two
pullbacks along explicit maps of pairs: on one side a point goes to the value
of the section at its chart image, on the other to the point of the fibre over
$z$ with the given coordinates.  These maps are homotopic through maps of
pairs, up to composing one of them with a real-linear automorphism of
$\C^{r} = \R^{2r}$ coming from the derivative of the section at $z$; the
automorphism is real-linear because the identification of the tangent space
with $\C^r$ comes from the real dimension count $6 + 2m = 2r$.  The relative
cohomology groups involved, that of the punctured neighbourhood and that of
$(\C^r, \C^r \setminus 0)$, are one-dimensional over $\F_2$, so any two
isomorphisms between them coincide.  Hence the homotopy can be replaced by
its two ends and the automorphism induces the identity, and the two pullbacks
agree.  The second statement follows from Lemma~\ref{lem:hres}.
\end{proof}

The last step is specific to $\F_2$: over any other field two isomorphisms
between lines differ by a scalar.

\section*{Provenance}

The problem was solved by Astra, using a custom proof tool.  The prose of
this paper was written by Claude Fable 5.1 (Anthropic).  The Lean
formalisation was carried out by AI agents under the author's direction.

\begin{thebibliography}{STW}
\raggedright

\bibitem{AJT}
S. Araki, I. M. James and E. Thomas,
Homotopy-abelian Lie groups,
\emph{Bull. Amer. Math. Soc.} \textbf{66} (1960), 324--326.

\bibitem{Bl}
B. Blackadar,
\emph{K-Theory for Operator Algebras}, 2nd ed.,
Math. Sci. Res. Inst. Publ. 5, Cambridge University Press, 1998.

\bibitem{Bo}
R. Bott,
The stable homotopy of the classical groups,
\emph{Ann. of Math. (2)} \textbf{70} (1959), 313--337.

\bibitem{Cu}
J. Cuntz,
K-theory for certain C*-algebras,
\emph{Ann. of Math. (2)} \textbf{113} (1981), 181--197.

\bibitem{EHT}
G. A. Elliott, T. M. Ho and A. S. Toms,
A class of simple C*-algebras with stable rank one,
\emph{J. Funct. Anal.} \textbf{256} (2009), 307--322.

\bibitem{Ji}
X. Jiang,
Nonstable K-theory for $\mathcal Z$-stable C*-algebras,
arXiv:math/9707228 (1997).

\bibitem{Li}
H. Lin,
Approximation by normal elements with finite spectra in C*-algebras of real
rank zero,
\emph{Pacific J. Math.} \textbf{173} (1996), 443--489.

\bibitem{MS}
J. W. Milnor and J. D. Stasheff,
\emph{Characteristic Classes},
Ann. of Math. Stud. 76, Princeton University Press, 1974.

\bibitem{Ri}
M. A. Rieffel,
The homotopy groups of the unitary groups of non-commutative tori,
\emph{J. Operator Theory} \textbf{17} (1987), 237--254.

\bibitem{STW}
C. Schafhauser, A. Tikuisis and S. White,
Nuclear C*-algebras: 99 problems,
arXiv:2506.10902 (2025).

\bibitem{Se}
A. Seth,
Tensorial permanence of $K$-stability for diagonal AH-algebras,
\emph{Bull. Lond. Math. Soc.} \textbf{58} (2026), \doi{10.1112/blms.70398}.

\bibitem{To}
A. S. Toms,
A simple C*-algebra which is not $K_1$-injective,
arXiv:2609.09535 (2026).

\bibitem{Vi98}
J. Villadsen,
Simple C*-algebras with perforation,
\emph{J. Funct. Anal.} \textbf{154} (1998), 110--116.

\bibitem{Vi99}
J. Villadsen,
On the stable rank of simple C*-algebras,
\emph{J. Amer. Math. Soc.} \textbf{12} (1999), 1091--1102.

\bibitem{Vi02}
J. Villadsen,
Comparison of projections and unitary elements in simple C*-algebras,
\emph{J. Reine Angew. Math.} \textbf{549} (2002), 23--45.

\end{thebibliography}

\end{document}
